\documentclass[14 pt]{extarticle}

	\usepackage[frenchb]{babel}
	\usepackage[utf8]{inputenc}  
	\usepackage[T1]{fontenc}
	\usepackage{amssymb}
	\usepackage[mathscr]{euscript}
	\usepackage{stmaryrd}
	\usepackage{amsmath}
	\usepackage{tikz}
	\usepackage[all,cmtip]{xy}
	\usepackage{amsthm}
	\usepackage{varioref}
	\usepackage{geometry}
	\geometry{a4paper}
	\usepackage{lmodern}
	\usepackage{hyperref}
	\usepackage{array}
	 \usepackage{fancyhdr}
	 \usepackage{float}
\renewcommand{\theenumi}{\alph{enumi})}
	\pagestyle{fancy}
	\theoremstyle{plain}
	\fancyfoot[C]{} 
	\fancyhead[L]{Contrôle}
	\fancyhead[R]{13 janvier 2025-2026}\geometry{
 a4paper,
 total={170mm,257mm},
 left=20mm,
 top=20mm,
 }
	
	
	\title{Interrogation chapitre 4}
	\date{}
	\begin{document}

\begin{center}{\Large Interrogation chapitre 4}\\ 
 \end{center}
 
 
 \subsection*{Exercice 1 (10 points)}
 
\begin{enumerate} 
\item Pour chacun des triplets de longueurs suivants, dire s'il permet de construire un triangle, en justifiant votre réponse.
\begin{enumerate}
\item $AB = 3$ cm, $AC = 7$ cm, $BC = 2$ cm ;
\item $AB = 3$ cm, $AC = 3$ cm, $BC = 5$ cm ;
\item $AB = 4$ cm, $AC = 3$ cm, $BC = 5$ cm. 
\end{enumerate}

\item Pour le ou les triangles dont vous avez montré l'existence, tracer le triangle, puis tracer ses médiatrices et son cercle circonscrit. 
\end{enumerate}


\subsection*{Exercice 2 (5 points)}
\begin{enumerate}
\item Calculer l'écriture décimale de $\frac2{11}$. 
\item Calculer $15 - 1 +4$
\item Calculer $15 \div 3\times 5$
\item Simplifier $\frac{315}{735}$ le plus possible. 
\end{enumerate}


\subsection*{Exercice 3 (5 points)}
On veut ici établir la propriété suivante : \og Si un quadrilatère $ABCD$ a ses quatre sommets sur un même cercle, alors les médiatrices de ses quatre côtés sont concourantes.\fg
\begin{enumerate}
\item Faire une figure de cette situation (où $ABCD$ n'est ni un rectangle, ni un carré).
\item On note $O$ le centre du cercle passant par $A$, $B$, $C$ et $D$. Montrer que la médiatrice de $[AB]$ contient le point $O$. 
\item En déduire que les médiatrices des quatre côtés se croisent au point $O$. 
\item Que représente le point $O$ pour chacun des triangles $ABC$, $BCD$, $CDA$, et $DAB$ ?
\end{enumerate}

\newpage 
\begin{center}{\Large Interrogation chapitre 4}\\ 
 \end{center}
 
 
 \subsection*{Exercice 1 (10 points)} 
 
 \begin{enumerate}
 \item Pour chacun des triplets de longueurs suivants, dire s'il permet de construire un triangle, en justifiant votre réponse.  
\begin{enumerate}
\item $AB = 3$ cm, $AC = 4$ cm, $BC = 2$ cm ;
\item $AB = 3$ cm, $AC = 3$ cm, $BC = 7$ cm ;
\item $AB = 2,5$ cm, $AC = 7,5$ cm, $BC = 6$ cm. 
\end{enumerate}

\item Pour le ou les triangles dont vous avez montré l'existence, tracer le triangle, puis tracer ses médiatrices et son cercle circonscrit. 
\end{enumerate}


\subsection*{Exercice 2 (5 points)}
\begin{enumerate}
\item Calculer l'écriture décimale de $\frac2{11}$. 
\item Calculer $11 - 1 +4$
\item Calculer $30 \div 3\times 10$
\item Simplifier $\frac{525}{420}$ le plus possible. 
\end{enumerate}


\subsection*{Exercice 3 (5 points)}
On veut ici établir la propriété suivante : \og Si un quadrilatère $ABCD$ a ses quatre sommets sur un même cercle, alors les médiatrices de ses quatre côtés sont concourantes.\fg
\begin{enumerate}
\item Faire une figure de cette situation (où $ABCD$ n'est ni un rectangle, ni un carré).
\item On note $O$ le centre du cercle passant par $A$, $B$, $C$ et $D$. Montrer que la médiatrice de $[AB]$ contient le point $O$. 
\item En déduire que les médiatrices des quatre côtés se croisent au point $O$. 
\item Que représente le point $O$ pour chacun des triangles $ABC$, $BCD$, $CDA$, et $DAB$ ?
\end{enumerate}

 	\end{document}
